\section{家庭、劳作与知识的十五世纪证据}

本组材料把视线从朝廷决策移向家庭、劳作、环境与知识如何在文书中留下痕迹。土地、户籍、税役、河工、学校、宗教空间与地方灾情，往往只在征解、工程、诉讼、捐输或褒奖的片段中出现。材料的缺口并非无关紧要：它说明哪些人的行动更容易进入档案，哪些人的经验只能从地方记录和物质遗存中间接追索。

\subsection{家庭责任与征解}

\paragraph{1432（宣德七年）：明朝派遣一失哈向明朝联军女真人进贡印玺并修复永宁寺} 这一锚点使家庭、劳动或地方资源进入可追查的历史时间。账本注明的把握范围是编年候选的年序可由官修与后出材料交叉；具体月日、参与者、战果或数字必须回查所列材料，不能将单一叙事当作定论。《宣宗实录》宣德七年条；《明史·本纪》及相关列传；诏令、奏疏或会典版本 提供了定位事件的材料链，却分别反映编纂者、报送者和保存者的视角。其发生条件可概括为继承、官僚协调或皇权运作的压力使问题进入朝廷程序；随后可见的制度或社会影响是它影响后续人事与政策议程；动机和派系标签不能由单一叙事裁定。对于日常生活、性别与家庭分工、非精英劳作或未留名者，仍不能由这些材料直接补足。译写候选以编年材料定位；实录记载反映朝廷选择，崇祯期须另以奏疏、地方及敌对政权材料交叉。

\paragraph{1432（宣德七年）：朝贡、册封和边地交涉的记录为本年多政权关系留档，礼仪语言与实际控制范围必须区分} 这一锚点使家庭、劳动或地方资源进入可追查的历史时间。账本注明的把握范围是来源导向的年度桥接条目：本年材料可定位相关政务线索；具体月日、发文机关、地域和执行结果仍须逐项回查，不能以制度文本或奏报替代实际效果。《宣宗实录》宣德七年条；《明史·外国传》；《朝鲜王朝实录》、琉球《历代宝案》或海外档案 提供了定位事件的材料链，却分别反映编纂者、报送者和保存者的视角。其发生条件可概括为朝贡名义、边境安全与港口贸易的利益使交涉持续发生；随后可见的制度或社会影响是它为后续册封、互市或海防安排留下文书与跨政权比较线索。对于日常生活、性别与家庭分工、非精英劳作或未留名者，仍不能由这些材料直接补足。桥接条目只标示可回查的年度问题与材料链；实录有中央选择性，崇祯期尤其需要奏疏、地方、朝鲜及满文材料交叉。

\paragraph{1432（宣德七年）：漕运、粮饷与军户事务的年度文书显示财政—军事相互约束，定额不能替代实际征解} 这一锚点使家庭、劳动或地方资源进入可追查的历史时间。账本注明的把握范围是来源导向的年度桥接条目：本年材料可定位相关政务线索；具体月日、发文机关、地域和执行结果仍须逐项回查，不能以制度文本或奏报替代实际效果。《宣宗实录》宣德七年条；《明会典》相关条目；地方志、黄册/契约/河工材料 提供了定位事件的材料链，却分别反映编纂者、报送者和保存者的视角。其发生条件可概括为财政征解、交通工程或货币支付的压力使相关措施进入政务；随后可见的制度或社会影响是其法定安排不等于地方执行，后续须以地域材料追踪负担与成效。对于日常生活、性别与家庭分工、非精英劳作或未留名者，仍不能由这些材料直接补足。桥接条目只标示可回查的年度问题与材料链；实录有中央选择性，崇祯期尤其需要奏疏、地方、朝鲜及满文材料交叉。

\paragraph{1433（宣德八年）：宝藏航行：郑和去世} 这一锚点使家庭、劳动或地方资源进入可追查的历史时间。账本注明的把握范围是编年候选的年序可由官修与后出材料交叉；具体月日、参与者、战果或数字必须回查所列材料，不能将单一叙事当作定论。《宣宗实录》宣德八年条；《明史·本纪》及相关列传；诏令、奏疏或会典版本 提供了定位事件的材料链，却分别反映编纂者、报送者和保存者的视角。其发生条件可概括为继承、官僚协调或皇权运作的压力使问题进入朝廷程序；随后可见的制度或社会影响是它影响后续人事与政策议程；动机和派系标签不能由单一叙事裁定。对于日常生活、性别与家庭分工、非精英劳作或未留名者，仍不能由这些材料直接补足。译写候选以编年材料定位；实录记载反映朝廷选择，崇祯期须另以奏疏、地方及敌对政权材料交叉。

\paragraph{1433（宣德八年）：宝藏之旅：洪宝和马欢抵达卡利卡特并派七人前往麦加，洪宝访问了乔法尔、拉萨、亚丁、摩加迪沙和巴拉瓦，然后返回中国} 这一锚点使家庭、劳动或地方资源进入可追查的历史时间。账本注明的把握范围是编年候选的年序可由官修与后出材料交叉；具体月日、参与者、战果或数字必须回查所列材料，不能将单一叙事当作定论。《宣宗实录》宣德八年条；《明史·本纪》及相关列传；诏令、奏疏或会典版本 提供了定位事件的材料链，却分别反映编纂者、报送者和保存者的视角。其发生条件可概括为继承、官僚协调或皇权运作的压力使问题进入朝廷程序；随后可见的制度或社会影响是它影响后续人事与政策议程；动机和派系标签不能由单一叙事裁定。对于日常生活、性别与家庭分工、非精英劳作或未留名者，仍不能由这些材料直接补足。译写候选以编年材料定位；实录记载反映朝廷选择，崇祯期须另以奏疏、地方及敌对政权材料交叉。

\paragraph{1433（宣德八年）三月9日：宝藏之旅：宝船队从霍尔木兹出发返回中国} 这一锚点使家庭、劳动或地方资源进入可追查的历史时间。账本注明的把握范围是编年候选的年序可由官修与后出材料交叉；具体月日、参与者、战果或数字必须回查所列材料，不能将单一叙事当作定论。《宣宗实录》宣德八年条；《明史·本纪》及相关列传；诏令、奏疏或会典版本 提供了定位事件的材料链，却分别反映编纂者、报送者和保存者的视角。其发生条件可概括为继承、官僚协调或皇权运作的压力使问题进入朝廷程序；随后可见的制度或社会影响是它影响后续人事与政策议程；动机和派系标签不能由单一叙事裁定。对于日常生活、性别与家庭分工、非精英劳作或未留名者，仍不能由这些材料直接补足。译写候选以编年材料定位；实录记载反映朝廷选择，崇祯期须另以奏疏、地方及敌对政权材料交叉。

\subsection{军役、道路与迁徙}

\paragraph{1433（宣德八年）六月：日本使团访问明朝中国：明日关系恢复} 这一锚点使家庭、劳动或地方资源进入可追查的历史时间。账本注明的把握范围是编年候选的年序可由官修与后出材料交叉；具体月日、参与者、战果或数字必须回查所列材料，不能将单一叙事当作定论。《宣宗实录》宣德八年条；《明史·外国传》；《朝鲜王朝实录》、琉球《历代宝案》或海外档案 提供了定位事件的材料链，却分别反映编纂者、报送者和保存者的视角。其发生条件可概括为朝贡名义、边境安全与港口贸易的利益使交涉持续发生；随后可见的制度或社会影响是它为后续册封、互市或海防安排留下文书与跨政权比较线索。对于日常生活、性别与家庭分工、非精英劳作或未留名者，仍不能由这些材料直接补足。译写候选以编年材料定位；实录记载反映朝廷选择，崇祯期须另以奏疏、地方及敌对政权材料交叉。

\paragraph{1433（宣德八年）七月7日：宝藏之旅：宝船队回到中国} 这一锚点使家庭、劳动或地方资源进入可追查的历史时间。账本注明的把握范围是编年候选的年序可由官修与后出材料交叉；具体月日、参与者、战果或数字必须回查所列材料，不能将单一叙事当作定论。《宣宗实录》宣德八年条；《明史·本纪》及相关列传；诏令、奏疏或会典版本 提供了定位事件的材料链，却分别反映编纂者、报送者和保存者的视角。其发生条件可概括为继承、官僚协调或皇权运作的压力使问题进入朝廷程序；随后可见的制度或社会影响是它影响后续人事与政策议程；动机和派系标签不能由单一叙事裁定。对于日常生活、性别与家庭分工、非精英劳作或未留名者，仍不能由这些材料直接补足。译写候选以编年材料定位；实录记载反映朝廷选择，崇祯期须另以奏疏、地方及敌对政权材料交叉。

\paragraph{1433（宣德八年）九月14日：宝藏航行：萨穆德拉·帕赛苏丹国、卡利卡特、科钦、锡兰、乔法尔、亚丁、哥印拜陀、霍尔木兹、卡亚尔和麦加的使节进贡} 这一锚点使家庭、劳动或地方资源进入可追查的历史时间。账本注明的把握范围是编年候选的年序可由官修与后出材料交叉；具体月日、参与者、战果或数字必须回查所列材料，不能将单一叙事当作定论。《宣宗实录》宣德八年条；《明史·外国传》；《朝鲜王朝实录》、琉球《历代宝案》或海外档案 提供了定位事件的材料链，却分别反映编纂者、报送者和保存者的视角。其发生条件可概括为朝贡名义、边境安全与港口贸易的利益使交涉持续发生；随后可见的制度或社会影响是它为后续册封、互市或海防安排留下文书与跨政权比较线索。对于日常生活、性别与家庭分工、非精英劳作或未留名者，仍不能由这些材料直接补足。译写候选以编年材料定位；实录记载反映朝廷选择，崇祯期须另以奏疏、地方及敌对政权材料交叉。

\paragraph{1433（宣德八年）：宝藏之旅：马欢出版《莺芽圣澜》} 这一锚点使家庭、劳动或地方资源进入可追查的历史时间。账本注明的把握范围是编年候选的年序可由官修与后出材料交叉；具体月日、参与者、战果或数字必须回查所列材料，不能将单一叙事当作定论。《宣宗实录》宣德八年条；《明史·本纪》及相关列传；诏令、奏疏或会典版本 提供了定位事件的材料链，却分别反映编纂者、报送者和保存者的视角。其发生条件可概括为继承、官僚协调或皇权运作的压力使问题进入朝廷程序；随后可见的制度或社会影响是它影响后续人事与政策议程；动机和派系标签不能由单一叙事裁定。对于日常生活、性别与家庭分工、非精英劳作或未留名者，仍不能由这些材料直接补足。译写候选以编年材料定位；实录记载反映朝廷选择，崇祯期须另以奏疏、地方及敌对政权材料交叉。

\paragraph{1434（宣德九年）：宝藏之旅：龚震出版《西洋繁国志》} 这一锚点使家庭、劳动或地方资源进入可追查的历史时间。账本注明的把握范围是编年候选的年序可由官修与后出材料交叉；具体月日、参与者、战果或数字必须回查所列材料，不能将单一叙事当作定论。《宣宗实录》宣德九年条；《明史·本纪》及相关列传；诏令、奏疏或会典版本 提供了定位事件的材料链，却分别反映编纂者、报送者和保存者的视角。其发生条件可概括为继承、官僚协调或皇权运作的压力使问题进入朝廷程序；随后可见的制度或社会影响是它影响后续人事与政策议程；动机和派系标签不能由单一叙事裁定。对于日常生活、性别与家庭分工、非精英劳作或未留名者，仍不能由这些材料直接补足。译写候选以编年材料定位；实录记载反映朝廷选择，崇祯期须另以奏疏、地方及敌对政权材料交叉。

\paragraph{1434（宣德九年）：本年朝廷围绕继承、官僚协调和政令传达形成连续记录，可用于核对宣德至正统的中枢运作} 这一锚点使家庭、劳动或地方资源进入可追查的历史时间。账本注明的把握范围是来源导向的年度桥接条目：本年材料可定位相关政务线索；具体月日、发文机关、地域和执行结果仍须逐项回查，不能以制度文本或奏报替代实际效果。《宣宗实录》宣德九年条；《明史·本纪》及相关列传；诏令、奏疏或会典版本 提供了定位事件的材料链，却分别反映编纂者、报送者和保存者的视角。其发生条件可概括为继承、官僚协调或皇权运作的压力使问题进入朝廷程序；随后可见的制度或社会影响是它影响后续人事与政策议程；动机和派系标签不能由单一叙事裁定。对于日常生活、性别与家庭分工、非精英劳作或未留名者，仍不能由这些材料直接补足。桥接条目只标示可回查的年度问题与材料链；实录有中央选择性，崇祯期尤其需要奏疏、地方、朝鲜及满文材料交叉。

\subsection{灾害、工程与地方记忆}

\paragraph{1434（宣德九年）：学校、科举和礼制的本年政策材料为社会流动和国家知识秩序提供可检索锚点} 这一锚点使家庭、劳动或地方资源进入可追查的历史时间。账本注明的把握范围是来源导向的年度桥接条目：本年材料可定位相关政务线索；具体月日、发文机关、地域和执行结果仍须逐项回查，不能以制度文本或奏报替代实际效果。《宣宗实录》宣德九年条；《明史·选举志》或《艺文志》；文集、碑刻或书目材料 提供了定位事件的材料链，却分别反映编纂者、报送者和保存者的视角。其发生条件可概括为制度、教育或知识传播的需要推动了相关行动或文本形成；随后可见的制度或社会影响是它提供制度和传播史的锚点，不能据此推断社会整体接受。对于日常生活、性别与家庭分工、非精英劳作或未留名者，仍不能由这些材料直接补足。桥接条目只标示可回查的年度问题与材料链；实录有中央选择性，崇祯期尤其需要奏疏、地方、朝鲜及满文材料交叉。

\paragraph{1434（宣德九年）：朝贡、册封和边地交涉的记录为本年多政权关系留档，礼仪语言与实际控制范围必须区分} 这一锚点使家庭、劳动或地方资源进入可追查的历史时间。账本注明的把握范围是来源导向的年度桥接条目：本年材料可定位相关政务线索；具体月日、发文机关、地域和执行结果仍须逐项回查，不能以制度文本或奏报替代实际效果。《宣宗实录》宣德九年条；《明史·外国传》；《朝鲜王朝实录》、琉球《历代宝案》或海外档案 提供了定位事件的材料链，却分别反映编纂者、报送者和保存者的视角。其发生条件可概括为朝贡名义、边境安全与港口贸易的利益使交涉持续发生；随后可见的制度或社会影响是它为后续册封、互市或海防安排留下文书与跨政权比较线索。对于日常生活、性别与家庭分工、非精英劳作或未留名者，仍不能由这些材料直接补足。桥接条目只标示可回查的年度问题与材料链；实录有中央选择性，崇祯期尤其需要奏疏、地方、朝鲜及满文材料交叉。

\paragraph{1435（宣德十年）一月31日：宣德皇帝驾崩，张皇后（洪熙）摄政为正统皇帝} 这一锚点使家庭、劳动或地方资源进入可追查的历史时间。账本注明的把握范围是编年候选的年序可由官修与后出材料交叉；具体月日、参与者、战果或数字必须回查所列材料，不能将单一叙事当作定论。《宣宗实录》宣德十年条；《明史·本纪》及相关列传；诏令、奏疏或会典版本 提供了定位事件的材料链，却分别反映编纂者、报送者和保存者的视角。其发生条件可概括为继承、官僚协调或皇权运作的压力使问题进入朝廷程序；随后可见的制度或社会影响是它影响后续人事与政策议程；动机和派系标签不能由单一叙事裁定。对于日常生活、性别与家庭分工、非精英劳作或未留名者，仍不能由这些材料直接补足。译写候选以编年材料定位；实录记载反映朝廷选择，崇祯期须另以奏疏、地方及敌对政权材料交叉。

\paragraph{1435（宣德十年）：华北平原及山东遭遇旱灾、瘟疫} 这一锚点使家庭、劳动或地方资源进入可追查的历史时间。账本注明的把握范围是编年候选的年序可由官修与后出材料交叉；具体月日、参与者、战果或数字必须回查所列材料，不能将单一叙事当作定论。《宣宗实录》宣德十年条；《明史·五行志》；受灾府县地方志与奏疏 提供了定位事件的材料链，却分别反映编纂者、报送者和保存者的视角。其发生条件可概括为自然风险与地方报告促成朝廷或地方的应对记录；随后可见的制度或社会影响是赈济、迁徙和损失的实际范围仍须以受灾地区材料分别核验。对于日常生活、性别与家庭分工、非精英劳作或未留名者，仍不能由这些材料直接补足。译写候选以编年材料定位；实录记载反映朝廷选择，崇祯期须另以奏疏、地方及敌对政权材料交叉。

\paragraph{1435（宣德十年）：学校、科举和礼制的本年政策材料为社会流动和国家知识秩序提供可检索锚点} 这一锚点使家庭、劳动或地方资源进入可追查的历史时间。账本注明的把握范围是来源导向的年度桥接条目：本年材料可定位相关政务线索；具体月日、发文机关、地域和执行结果仍须逐项回查，不能以制度文本或奏报替代实际效果。《宣宗实录》宣德十年条；《明史·选举志》或《艺文志》；文集、碑刻或书目材料 提供了定位事件的材料链，却分别反映编纂者、报送者和保存者的视角。其发生条件可概括为制度、教育或知识传播的需要推动了相关行动或文本形成；随后可见的制度或社会影响是它提供制度和传播史的锚点，不能据此推断社会整体接受。对于日常生活、性别与家庭分工、非精英劳作或未留名者，仍不能由这些材料直接补足。桥接条目只标示可回查的年度问题与材料链；实录有中央选择性，崇祯期尤其需要奏疏、地方、朝鲜及满文材料交叉。

\paragraph{1436（正统元年）：宝藏航行：明朝禁止建造海船} 这一锚点使家庭、劳动或地方资源进入可追查的历史时间。账本注明的把握范围是编年候选的年序可由官修与后出材料交叉；具体月日、参与者、战果或数字必须回查所列材料，不能将单一叙事当作定论。《英宗实录》正统元年条；《明史·本纪》及相关列传；诏令、奏疏或会典版本 提供了定位事件的材料链，却分别反映编纂者、报送者和保存者的视角。其发生条件可概括为继承、官僚协调或皇权运作的压力使问题进入朝廷程序；随后可见的制度或社会影响是它影响后续人事与政策议程；动机和派系标签不能由单一叙事裁定。对于日常生活、性别与家庭分工、非精英劳作或未留名者，仍不能由这些材料直接补足。译写候选以编年材料定位；实录记载反映朝廷选择，崇祯期须另以奏疏、地方及敌对政权材料交叉。

\subsection{学校、文本与行政知识}

\paragraph{1436（正统元年）：宝藏之旅：费辛出版《星茶圣澜》} 这一锚点使家庭、劳动或地方资源进入可追查的历史时间。账本注明的把握范围是编年候选的年序可由官修与后出材料交叉；具体月日、参与者、战果或数字必须回查所列材料，不能将单一叙事当作定论。《英宗实录》正统元年条；《明史·本纪》及相关列传；诏令、奏疏或会典版本 提供了定位事件的材料链，却分别反映编纂者、报送者和保存者的视角。其发生条件可概括为继承、官僚协调或皇权运作的压力使问题进入朝廷程序；随后可见的制度或社会影响是它影响后续人事与政策议程；动机和派系标签不能由单一叙事裁定。对于日常生活、性别与家庭分工、非精英劳作或未留名者，仍不能由这些材料直接补足。译写候选以编年材料定位；实录记载反映朝廷选择，崇祯期须另以奏疏、地方及敌对政权材料交叉。

\paragraph{1436（正统元年）：苏北、华北平原、山东等地遭受洪涝灾害} 这一锚点使家庭、劳动或地方资源进入可追查的历史时间。账本注明的把握范围是编年候选的年序可由官修与后出材料交叉；具体月日、参与者、战果或数字必须回查所列材料，不能将单一叙事当作定论。《英宗实录》正统元年条；《明史·五行志》；受灾府县地方志与奏疏 提供了定位事件的材料链，却分别反映编纂者、报送者和保存者的视角。其发生条件可概括为自然风险与地方报告促成朝廷或地方的应对记录；随后可见的制度或社会影响是赈济、迁徙和损失的实际范围仍须以受灾地区材料分别核验。对于日常生活、性别与家庭分工、非精英劳作或未留名者，仍不能由这些材料直接补足。译写候选以编年材料定位；实录记载反映朝廷选择，崇祯期须另以奏疏、地方及敌对政权材料交叉。

\paragraph{1437（正统二年）：山西、陕西遭遇干旱} 这一锚点使家庭、劳动或地方资源进入可追查的历史时间。账本注明的把握范围是编年候选的年序可由官修与后出材料交叉；具体月日、参与者、战果或数字必须回查所列材料，不能将单一叙事当作定论。《英宗实录》正统二年条；《明史·五行志》；受灾府县地方志与奏疏 提供了定位事件的材料链，却分别反映编纂者、报送者和保存者的视角。其发生条件可概括为自然风险与地方报告促成朝廷或地方的应对记录；随后可见的制度或社会影响是赈济、迁徙和损失的实际范围仍须以受灾地区材料分别核验。对于日常生活、性别与家庭分工、非精英劳作或未留名者，仍不能由这些材料直接补足。译写候选以编年材料定位；实录记载反映朝廷选择，崇祯期须另以奏疏、地方及敌对政权材料交叉。

\paragraph{1437（正统二年）：苏北地区发生洪涝灾害} 这一锚点使家庭、劳动或地方资源进入可追查的历史时间。账本注明的把握范围是编年候选的年序可由官修与后出材料交叉；具体月日、参与者、战果或数字必须回查所列材料，不能将单一叙事当作定论。《英宗实录》正统二年条；《明史·五行志》；受灾府县地方志与奏疏 提供了定位事件的材料链，却分别反映编纂者、报送者和保存者的视角。其发生条件可概括为自然风险与地方报告促成朝廷或地方的应对记录；随后可见的制度或社会影响是赈济、迁徙和损失的实际范围仍须以受灾地区材料分别核验。对于日常生活、性别与家庭分工、非精英劳作或未留名者，仍不能由这些材料直接补足。译写候选以编年材料定位；实录记载反映朝廷选择，崇祯期须另以奏疏、地方及敌对政权材料交叉。

\paragraph{1437（正统二年）：灾情上报、赈恤或河工材料标记本年地方风险，损失与救济效果需回到受灾地区核实。（材料核查重点：诏令或奏议的形成与发受链）} 这一锚点使家庭、劳动或地方资源进入可追查的历史时间。账本注明的把握范围是来源导向的年度桥接条目：本年材料可定位相关政务线索；具体月日、发文机关、地域和执行结果仍须逐项回查，不能以制度文本或奏报替代实际效果。《英宗实录》正统二年条；《明史·五行志》；受灾府县地方志与奏疏 提供了定位事件的材料链，却分别反映编纂者、报送者和保存者的视角。其发生条件可概括为自然风险与地方报告促成朝廷或地方的应对记录；随后可见的制度或社会影响是赈济、迁徙和损失的实际范围仍须以受灾地区材料分别核验。对于日常生活、性别与家庭分工、非精英劳作或未留名者，仍不能由这些材料直接补足。桥接条目只标示可回查的年度问题与材料链；实录有中央选择性，崇祯期尤其需要奏疏、地方、朝鲜及满文材料交叉。；本条特别核对诏令或奏议的形成与发受链。

\paragraph{1437（正统二年）：灾情上报、赈恤或河工材料标记本年地方风险，损失与救济效果需回到受灾地区核实。（材料核查重点：报称信息的来源、抄传与行政分类）} 这一锚点使家庭、劳动或地方资源进入可追查的历史时间。账本注明的把握范围是来源导向的年度桥接条目：本年材料可定位相关政务线索；具体月日、发文机关、地域和执行结果仍须逐项回查，不能以制度文本或奏报替代实际效果。《英宗实录》正统二年条；《明史·五行志》；受灾府县地方志与奏疏 提供了定位事件的材料链，却分别反映编纂者、报送者和保存者的视角。其发生条件可概括为自然风险与地方报告促成朝廷或地方的应对记录；随后可见的制度或社会影响是赈济、迁徙和损失的实际范围仍须以受灾地区材料分别核验。对于日常生活、性别与家庭分工、非精英劳作或未留名者，仍不能由这些材料直接补足。桥接条目只标示可回查的年度问题与材料链；实录有中央选择性，崇祯期尤其需要奏疏、地方、朝鲜及满文材料交叉。；本条特别核对报称信息的来源、抄传与行政分类。
